%%%%%%%%%%%%%%%%%%%%%%%%%%%%%%%%%%%%%%%%%%%%%%%%%%%%%%%%%%%%%%%%%%%%%%%%%%% 
% 
% Generic template for TFC/TFM/TFG/Tesis
% 
% By:
%  + Javier Macías-Guarasa.
%    Departamento de Electrónica
%    Universidad de Alcalá
%  + Roberto Barra-Chicote.
%    Departamento de Ingeniería Electrónica
%    Universidad Politécnica de Madrid
% 
% By: + Javier Macías-Guarasa. Departamento de Electrónica Universidad de Alcalá + Roberto Barra-Chicote. Departamento de Ingeniería Electrónica Universidad Politécnica de Madrid
% 
% Based on original sources by Roberto Barra, Manuel Ocaña, Jesús Nuevo, Pedro Revenga, Fernando Herránz and Noelia Hernández. Thanks a lot to all of them, and to the many anonymous contributors found (thanks to google) that provided help in setting all this up.
% 
% See also the additionalContributors.txt file to check the name of additional contributors to this work.
% 
% If you think you can add pieces of relevant/useful examples, improvements, please contact us at (macias@depeca.uah.es)
% 
% You can freely use this template and please contribute with comments or suggestions!!!
% 
%%%%%%%%%%%%%%%%%%%%%%%%%%%%%%%%%%%%%%%%%%%%%%%%%%%%%%%%%%%%%%%%%%%%%%%%%%% 

\chapter{Introducción}
\label{cha:introduccion}
%
%\begin{FraseCelebre}
%  \begin{Frase}
%    Desocupado lector, sin juramento me podrás creer que quisiera que este libro [...] fuera el más hermoso, el más gallardo y más discreto que pudiera imaginarse\footnote{Tomado de ejemplos del proyecto \texis{}.}.
%  \end{Frase}
%  \begin{Fuente}
%    Miguel de Cervantes, Don Quijote de la Mancha
%  \end{Fuente}
%\end{FraseCelebre}


\section{Contexto y motivación}
\label{sec:contexto-motivacion}

El consumo energético se ha convertido en una preocupación
importante dentro del debate sobre la sostenibilidad digital. Durante las
últimas décadas, casi todos los esfuerzos por gastar menos energía se han
centrado en el hardware, donde se han producido grandes avances.

Un ejemplo es la fabricación de los procesadores. Al reducir el tamaño de los
transistores se pueden colocar muchos más en el mismo espacio de chip, de modo que cada nueva generación ofrece más
potencia y, a la vez, gasta menos energía en realizar el mismo trabajo.

A esto se suma la evolución de las arquitecturas de procesador. Los diseños de
tipo \emph{big.LITTLE} \cite{armbiglittle}, propios de la arquitectura ARM,
combinan núcleos potentes con núcleos de bajo consumo y reparten el trabajo
entre ellos según lo que haga falta en cada momento, de forma que las tareas
ligeras apenas consumen batería.

Este enfoque nació en el mundo de los dispositivos móviles, pero ha funcionado
tan bien que la arquitectura ARM está dando el salto a los ordenadores. El caso
más destacado fueron los chips M de Apple (a partir del M1, en 2020), que
llevaron esta filosofía a portátiles y sobremesas con un rendimiento y una
eficiencia que supusieron una revolución en el sector tecnológico. Desde entonces otros fabricantes han seguido el mismo camino, como Qualcomm, que con sus procesadores Snapdragon ha trasladado al PC la misma arquitectura que ya usaba en los móviles. Es más, la idea de combinar núcleos de alto rendimiento con núcleos de
bajo consumo ha calado incluso fuera de ARM, y fabricantes como Intel la han
incorporado a sus procesadores más recientes manteniéndose en la arquitectura
x86. Y más allá del ordenador destinado a uso personal, su eficiencia ha impulsado su
adopción también en servidores y centros de datos, donde reducir el consumo
supone un gran ahorro de costes. Un ejemplo claro es
Amazon, que diseñó sus propios procesadores ARM, conocidos como AWS Graviton,
para sus servidores en la nube, y que según la compañía llegan a consumir hasta
un 60\% menos de energía que instancias equivalentes con el mismo rendimiento
\cite{awsgraviton}.

El software, en cambio, casi nunca se diseña ni se evalúa pensando en la
energía que gasta, a pesar de que es él quien decide cómo se usan todos esos
recursos del hardware \cite{verdecchia2021}.

Y esto tiene consecuencias reales. Un mismo problema se puede resolver de
varias formas distintas que funcionan igual de bien, pero que consumen cantidades
de energía muy diferentes. Es decir, las decisiones que toma quien programa
afectan directamente a la energía que gasta el programa al ejecutarse
\cite{noureddine2025}. A pequeña escala, esto se nota en la batería de nuestros dispositivos, y a gran escala, en el consumo eléctrico de los grandes centros de datos, que no para de crecer.

Para poder mejorar esto hace falta, antes de nada, poder medirlo. Así nace la
llamada \emph{ingeniería del software sostenible} o \emph{green software}, que
busca tratar la eficiencia energética como un criterio de calidad más, igual
que lo son el rendimiento o la seguridad. Y para eso resultan imprescindibles
herramientas capaces de medir cuánta energía consume una aplicación de forma
fiable y con detalle, de manera que las decisiones de diseño se apoyen en datos
y no en suposiciones. Este Trabajo Fin de Grado se enmarca precisamente en esa
necesidad.



\section{Descripción del problema}
\label{sec:descripcion-problema}

La cantidad de energía que consume un programa no es fácil de medir. Existen
varias formas de conocer el consumo del equipo en su conjunto, como usar
medidores físicos que se conectan a la corriente, o programas que leen los
sensores internos del ordenador (como HWiNFO o HWMonitor) y estiman los vatios
que gasta el equipo en tiempo real. El problema de estas herramientas es que miden el consumo de todo el equipo, por lo que no es posible saber exactamente cuál es el de una aplicación o programa en concreto que se está ejecutando. Además, un programa está formado
internamente por muchas partes o funciones que se encargan de realizar tareas
distintas (como leer un fichero o procesar datos), y cada una de ellas gasta
una cantidad de energía diferente.



Por tanto, para mejorar la eficiencia energética de un software, no basta con
saber cuánto consume el equipo a nivel global, sino que habría que poder medir
cada programa por separado y conocer qué partes son las que más consumen, para
así optimizarlas. Los métodos anteriores no son suficientes para ello, y es aquí
donde aparecen herramientas más específicas para este problema, como JoularJX.
JoularJX es un agente Java de código abierto desarrollado por el investigador
Adel Noureddine en la Université de Pau et des Pays de l'Adour
\cite{joularjx2022, joularjxrepo}, el cual permite medir,
de forma aproximada, cuánta energía consume cada función de un programa por
separado, y así identificar cuáles son las que más gastan para poder optimizarlas \cite{joularjx3}.

De hecho, que JoularJX esté diseñado específicamente para Java es un detalle importante, ya que es uno de los lenguajes más utilizados a nivel global (sobre todo en servidores y aplicaciones
empresariales que están en uso durante todo el día), por lo que cualquier mejora en
su consumo, por pequeña que sea, se multiplica por miles de equipos, lo que a gran escala supone un ahorro considerable. 

Además de estudiar JoularJX en sí, se aplicará a un caso real, comparando el
consumo energético de \textbf{Jackson} \cite{jackson} y \textbf{Gson}
\cite{gson}, dos de las librerías de Java más conocidas para procesar datos en
formato JSON, que es el formato más habitual para intercambiar datos entre aplicaciones. Las dos hacen lo mismo, pero están programadas de forma muy
distinta, lo que permitirá comparar su consumo y ver las diferencias entre ambas
\cite{noureddine2025}.

\section{Objetivos}
\label{sec:objetivos}


El objetivo de este trabajo es valorar hasta qué punto JoularJX es útil y fiable
para medir el consumo energético de aplicaciones Java, aplicándolo a la
comparación de Jackson y Gson.

Para alcanzar esta meta, se definen los siguientes objetivos específicos:

\begin{itemize}
	\item \textbf{Entender cómo mide JoularJX la energía que utiliza cada función}
	de un programa Java, viendo cómo funciona como agente, cómo detecta qué parte
	del programa se está ejecutando en cada momento y de dónde obtiene los datos
	de energía.
	\item \textbf{Instalar y configurar JoularJX en un entorno GNU/Linux},
	documentando el proceso paso a paso y comprobando que el sistema puede acceder a los datos de consumo del procesador.
	\item \textbf{Preparar un banco de pruebas con Jackson y Gson} para medir su consumo en operaciones equivalentes de
	lectura y escritura de datos JSON.
	\item \textbf{Ejecutar las mediciones y comparar ambas librerías}, analizando
	el consumo total, el de cada función y qué partes gastan más energía.
	\item \textbf{Evaluar las limitaciones de JoularJX}, como la precisión de sus
	mediciones, el consumo extra que añade la propia herramienta al utilizarse (ya que se ejecuta dentro del mismo programa y podría alterar los resultados) o cuánto cambian los valores obtenidos al repetir una misma prueba varias veces.
\end{itemize}


\section{Estructura de la memoria}
\label{sec:estructura}

A continuación se resume brevemente el contenido de cada uno de los capítulos que componen esta memoria:

\begin{itemize}
	\item El \textbf{Capítulo~\ref{cha:estudio-teorico}} explica en detalle el
	funcionamiento interno de JoularJX y los conceptos necesarios para
	entender cómo estima el consumo energético.
	\item El \textbf{Capítulo~\ref{cha:desarrollo}} describe la parte práctica:
	la instalación de la herramienta y la preparación del banco de pruebas con
	Jackson y Gson.
	\item El \textbf{Capítulo~\ref{cha:resultados}} presenta las mediciones
	obtenidas y la comparación del consumo entre ambas librerías.
	\item El \textbf{Capítulo~\ref{cha:concl-y-line}} recoge las conclusiones
	del trabajo, las limitaciones encontradas y las posibles mejoras o
	ampliaciones que se podrían hacer en el futuro.
\end{itemize}

%%% Local Variables:
%%% TeX-master: "../book"
%%% End:
